% =====================================================================
% Conteudo das secoes para colar no documento da faculdade
% (sem preambulo, sem capa, sem \begin{document})
%
% Pre-requisitos no preambulo do .tex principal:
%   \usepackage{enumitem}
%   \usepackage{booktabs}
% =====================================================================

\section{Mutex e Semáforo em Java}

A linguagem Java oferece, no pacote \texttt{java.util.concurrent}, as
classes \texttt{ReentrantLock} e \texttt{Semaphore}, que correspondem
diretamente ao \texttt{pthread\_mutex\_t} e ao \texttt{sem\_t} de POSIX.
A correspondência entre as primitivas é direta:

\begin{center}
\small
\begin{tabular}{ll}
\toprule
\textbf{C / POSIX} & \textbf{Java / \texttt{java.util.concurrent}} \\
\midrule
\texttt{pthread\_mutex\_t lock;}        & \texttt{ReentrantLock lock = new ReentrantLock();} \\
\texttt{pthread\_mutex\_lock(\&lock);}  & \texttt{lock.lock();} \\
\texttt{pthread\_mutex\_unlock(\&lock);}& \texttt{lock.unlock();} \\
\texttt{sem\_t s; sem\_init(\&s,0,N);}  & \texttt{Semaphore s = new Semaphore(N);} \\
\texttt{sem\_wait(\&s);}                & \texttt{s.acquire();} \\
\texttt{sem\_post(\&s);}                & \texttt{s.release();} \\
\bottomrule
\end{tabular}
\end{center}

A diferença prática é que em Java o \texttt{unlock()} (ou o \texttt{release()})
deve ficar em um bloco \texttt{finally} --- assim o lock é liberado mesmo se
a seção crítica lançar exceção, evitando deadlocks por vazamento.

\subsection{Mutex com \texttt{ReentrantLock}}

Reescrita do \texttt{m8-mutex.c}: 100 \textit{threads} incrementam um
contador 100\,000 vezes cada uma. Sem proteção, várias atualizações são
perdidas; com \texttt{ReentrantLock} o resultado fecha em
$100\,000 \times 100 = 10\,000\,000$. O código (\texttt{MutexExemplo.java})
segue o padrão \texttt{lock.lock()} antes da seção crítica e
\texttt{lock.unlock()} dentro de um \texttt{finally}.

\subsection{Semáforo binário com \texttt{Semaphore(1)}}

Reescrita do \texttt{m5-semaforo.c}. Com uma única ``permissão'' o semáforo
funciona como um mutex: a primeira \textit{thread} a chamar
\texttt{s.acquire()} consome a vaga; as demais ficam bloqueadas até o
\texttt{s.release()} correspondente. O programa
(\texttt{SemaforoExemplo.java}) executa as mesmas 100 \textit{threads}
$\times$ 100\,000 incrementos do exemplo anterior e também alcança o
total esperado de 10\,000\,000.

\subsection{Semáforo contador com \texttt{Semaphore(3)}}

Reescrita do \texttt{m6-semaforo.c}. O valor inicial $N=3$ permite no
máximo três \textit{threads} simultaneamente dentro da seção crítica
--- útil para limitar acesso concorrente a um recurso (banco de dados,
\textit{pool} de conexões, vagas de estacionamento). Em
\texttt{BancoConexoes.java} criei o semáforo como
\texttt{new Semaphore(3, true)} (modo \textit{fair}), o que serve as
\textit{threads} em ordem de chegada e reduz \textit{starvation} ao custo
de um pouco de \textit{throughput}. Nos testes com 10 \textit{threads},
nunca houve mais que 3 conectadas ao mesmo tempo.

% =====================================================================
\section{Conta bancária}

O \texttt{conta.c} original tem uma condição de disputa: a função
\texttt{depositar} faz três passos (ler, somar, escrever) com um
\texttt{usleep(100)} entre eles, forçando a troca de contexto. Em Java
reproduzi a mesma ideia com \texttt{Thread.sleep(0, 100\_000)}
(100\,$\mu$s de pausa entre a leitura e a escrita do saldo). A seguir
comparo a versão sem proteção contra as três soluções pedidas.

\subsection{Sem proteção --- a disputa exposta}

\texttt{ContaProblema.java} reproduz o cenário do \texttt{conta.c}: duas
\textit{threads} depositam simultaneamente em uma mesma variável
\texttt{saldo}, sem nenhum mecanismo de sincronização. Como
\texttt{t1} lê \texttt{saldo=100}, dorme, e nesse intervalo \texttt{t2}
também lê 100 e grava 130, quando \texttt{t1} retoma sobrescreve com
$100+50=150$. O depósito de \texttt{t2} é perdido. Em todas as 5
execuções de teste o saldo final ficou em 130, evidenciando a perda do
depósito de 50. O problema é determinístico aqui porque o \texttt{sleep}
força a interleaving.

\subsection{Solução com Mutex (\texttt{ReentrantLock})}

\texttt{ContaMutex.java} envolve as três operações da função
\texttt{depositar} (leitura, soma, escrita) entre \texttt{lock.lock()} e
\texttt{lock.unlock()} (em \texttt{finally}). Com isso, mesmo que o
\texttt{sleep} continue forçando a troca de contexto, a outra
\textit{thread} fica bloqueada na entrada do método e o saldo final é
sempre 180.

\subsection{Solução com Semáforo}

\texttt{ContaSemaforo.java} usa um \texttt{Semaphore(1)}, equivalente
funcional ao mutex. \texttt{s.acquire()} antes da seção crítica e
\texttt{s.release()} dentro de um \texttt{finally}. Resultado idêntico ao
do \texttt{ReentrantLock}: 180 em todas as execuções.

\subsection{Solução com Monitor (\texttt{synchronized})}

Aqui aplico a dica do enunciado: \texttt{Thread.sleep(0)} dentro do método
\texttt{synchronized}. Em \texttt{ContaMonitor.java} a conta é
encapsulada em uma classe interna e o método \texttt{depositar} é
\texttt{synchronized}. Como o lock do monitor \textbf{não} é liberado
durante o \texttt{sleep} (diferente de \texttt{wait()}), a outra
\textit{thread} fica bloqueada na entrada do método e o resultado segue
correto (180).

% =====================================================================
\section{Pizzaria com buffer de 2 pizzas}

A \texttt{Mesa} original (\texttt{Mesa.java}) só comporta uma pizza por
vez --- enquanto o cozinheiro não retira a anterior, ele bloqueia. Para
permitir que o cozinheiro ``estoque'' até duas pizzas, transformei a
mesa em um buffer com fila usando \texttt{Queue<String>} e capacidade
\texttt{CAP=2}, no arquivo \texttt{MesaBuffer.java}. Os métodos pedidos
pelo enunciado são todos usados:

\begin{itemize}[noitemsep, topsep=2pt, leftmargin=*]
  \item \texttt{fila.add(pizza)} --- cozinheiro adiciona uma pizza nova;
  \item \texttt{fila.size()} --- usado para checar se atingiu a capacidade máxima;
  \item \texttt{fila.isEmpty()} --- entregador checa se há algo a retirar;
  \item \texttt{fila.poll()} --- entregador retira a pizza mais antiga (FIFO).
\end{itemize}

A sincronização passou a usar \texttt{notifyAll()} no lugar de
\texttt{notify()}, garantindo que qualquer \textit{thread} bloqueada em
\texttt{wait()} (cozinheiro esperando vaga, ou entregador esperando
pizza) seja acordada quando o estado muda.

O driver \texttt{PizzariaBuffer.java} aceita a flag
\texttt{-DcozinheiroLento=true|false}, que troca os tempos de
\texttt{Thread.sleep} para testar os dois extremos.

\subsection{Cenário 1: cozinheiro rápido, entregador lento}

Tempo de preparo $= 200$\,ms; tempo de entrega $= 1500$\,ms. A fila
enche até $2/2$ e o cozinheiro precisa esperar o entregador retirar uma
pizza --- mensagens ``mesa cheia (2/2)'' aparecem repetidamente no log.

\subsection{Cenário 2: cozinheiro lento, entregador rápido}

Tempo de preparo $= 1500$\,ms; tempo de entrega $= 200$\,ms. A fila quase
sempre fica vazia: o entregador retira a única pizza e fica em
\texttt{wait()} até o cozinheiro produzir a próxima --- mensagens
``mesa vazia'' aparecem entre cada pizza.

\subsection{Por que \texttt{notifyAll()} em vez de \texttt{notify()}?}

Com um único produtor e um único consumidor, \texttt{notify()} bastaria.
Mas o exemplo é facilmente generalizável para vários cozinheiros e
entregadores, e nesse caso \texttt{notify()} pode acordar a
\textit{thread} ``errada'' (ex.: acordar outro produtor quando quem
estava esperando era um consumidor), causando deadlock.
\texttt{notifyAll()} acorda todas as \textit{threads} bloqueadas no
monitor; o \texttt{while} em volta do \texttt{wait()} cuida de checar a
condição novamente após o \textit{wake-up}.

% =====================================================================
\section{Limitações dos semáforos}

Apesar de teoricamente suficientes para qualquer problema de
sincronização, semáforos são primitivas de baixo nível com armadilhas
conhecidas. A seguir listo as principais limitações encontradas em
Tanenbaum (\textit{Modern Operating Systems}) e Silberschatz
(\textit{Operating System Concepts}), e o que cada uma pode provocar.

\begin{enumerate}[leftmargin=*, itemsep=4pt]

\item \textbf{Não há associação explícita entre semáforo e recurso.}
Um semáforo é apenas um inteiro com \texttt{wait}/\texttt{signal}.
O compilador não verifica se o programador adquiriu o semáforo antes
de tocar no dado protegido. Foi exatamente essa limitação que motivou
o conceito de Monitor (Hoare, 1974), que encapsula dado e seção crítica.

\item \textbf{Inversão / esquecimento de \texttt{wait} ou \texttt{signal}.}
Esquecer um \texttt{wait} expõe a região crítica (race silenciosa);
esquecer um \texttt{signal} causa deadlock; trocar a ordem causa
comportamento errático.

\item \textbf{Deadlock por ordem de aquisição.}
Quando uma tarefa precisa de mais de um semáforo, a ordem importa.
O jantar dos filósofos é o exemplo clássico: cada um pega o garfo da
esquerda e tenta o da direita; todos travam.

\item \textbf{Starvation.}
Sem fila justa (FIFO), uma \textit{thread} pode nunca ser acordada
quando há disputa. Em Java mitiga-se com
\texttt{new Semaphore(n, true)}, ao custo de \textit{throughput}.

\item \textbf{Inversão de prioridade.}
Uma \textit{thread} de baixa prioridade que detém o semáforo bloqueia
uma de alta prioridade. Caso famoso: Mars Pathfinder (1997). A solução
exige protocolos como \textit{priority inheritance}, externos ao
semáforo padrão.

\item \textbf{Composição difícil.}
Operações compostas (ex.: ``transferir entre duas contas atomicamente'')
não saem naturais com semáforos. Locks reentrantes e monitores compõem
melhor.

\item \textbf{Não-reentrância.}
Semáforos POSIX não são reentrantes: a mesma \textit{thread}, ao tentar
dar um segundo \texttt{wait} em um semáforo binário que já detém, trava
em si mesma. (O \texttt{Semaphore} de Java segue o mesmo princípio;
quem precisa de reentrância usa \texttt{ReentrantLock} ou
\texttt{synchronized}.)

\item \textbf{Espera ocupada em implementações ingênuas.}
Implementações didáticas em \textit{user-space} viram \texttt{while
(s==0);}, queimando CPU. Implementações reais usam suporte do SO, mas
o conceito ainda demanda cuidado em código de baixo nível.

\item \textbf{Ausência de espera por condição complexa.}
Semáforos sinalizam apenas ``há vaga / não há vaga''. Para condições
como ``buffer com pelo menos 3 itens \textit{e} horário antes das 18h'',
o programador precisa combinar vários semáforos --- propenso a erros.
Variáveis de condição (\texttt{wait}/\texttt{notify}) e monitores
resolvem com mais clareza.

\end{enumerate}

\subsection{Resumo das consequências}

\begin{center}
\small
\begin{tabular}{ll}
\toprule
\textbf{Falha de uso}                  & \textbf{Consequência}                  \\
\midrule
Esquecer \texttt{wait}                 & Race condition, dados corrompidos      \\
Esquecer \texttt{signal}               & Deadlock                               \\
Ordem inversa de locks                 & Deadlock circular                      \\
Sem \textit{fairness}                  & Starvation                             \\
Prioridades sem protocolo              & Inversão de prioridade                 \\
Recurso não encapsulado                & Bug invisível para o compilador        \\
\bottomrule
\end{tabular}
\end{center}

% =====================================================================
\section{Conclusão}

A reescrita em Java mostrou que \texttt{ReentrantLock} e \texttt{Semaphore}
oferecem o mesmo poder das primitivas POSIX usadas em C, com a vantagem
do \texttt{try/finally} (que evita deadlocks por exceção). No problema da
conta bancária, os três mecanismos --- mutex, semáforo e monitor ---
eliminaram a condição de disputa e produziram saldo final 180 em todas as
execuções, sendo o monitor a opção mais idiomática em Java.

A versão da pizzaria com buffer evidencia a vantagem do produtor/consumidor
desacoplado por uma fila: o cozinheiro não fica bloqueado em cada pizza,
podendo estocar até duas. Os dois cenários de teste (\textit{cozinheiro
lento} vs.\ \textit{rápido}) deixaram claro o comportamento esperado de
cada lado bloqueando quando o buffer fica vazio ou cheio.

Por fim, ainda que semáforos resolvam, em teoria, qualquer problema de
sincronização, suas limitações práticas (deadlock, starvation, inversão de
prioridade, falta de encapsulamento) explicam por que as linguagens
modernas oferecem mecanismos de mais alto nível como
\texttt{synchronized}, \texttt{BlockingQueue} e \texttt{ReentrantLock}.
